\cgeentry{cge_options}{Settings of the solver}

\cgesec{Description}

\cmd{cge_options} (from version 0.10.1) sets how the model is solved, as
GAMS's \cmd{option} statement chooses a solver; without options it displays
the settings. They are saved with the model and apply to every solve:
\manref{cge_solve}, \manref{cge_dynamic}, \manref{cge_sensitivity},
\manref{cge_optimize} and \manref{cge_loop}. \opt{lin}{solver()} chooses the
linear solver of Newton's steps: Mata's dense LU, the default, or a sparse LU,
whose memory grows with the number of entries of the Jacobian rather than with
the square of the number of equations. From version 0.14.0,
\opt{}{mcp()} chooses the method for constraints and inequalities: the
semismooth Newton method, the default, or the active-set method of earlier
versions.

\cgesec{Quick start}

The sparse LU for every later solve
\begin{stcode}
cge_options, linsolver(sparse)
\end{stcode}
The sparse LU for systems of 1,000 equations or more
\begin{stcode}
cge_options, linsolver(auto)
\end{stcode}
The active-set method for constraints
\begin{stcode}
cge_options, mcp(activeset)
\end{stcode}
Display the settings
\begin{stcode}
cge_options
\end{stcode}

\cgesec{Syntax}

\begin{stsyntax}
\cmdb{cge_options} [\texttt{,} \optarg{lin}{solver}{dense|sparse|auto} \optarg{}{mcp}{semismooth|activeset}]
\end{stsyntax}

\begin{opttable}
\optgroup{Main}
\optarg{lin}{solver}{dense|sparse|auto} & the linear solver of Newton's steps\\
\optarg{}{mcp}{semismooth|activeset} & the method for constraints and inequalities\\
\end{opttable}

\cgesec{Options}

\begin{optdesc}
\item[\optarg{lin}{solver}{dense|sparse|auto}] chooses the linear solver of
Newton's steps: \cmd{dense}, Mata's LU factorisation with partial pivoting
(the default); \cmd{sparse}, the sparse LU of section~\ref{sec:sparse}; or
\cmd{auto}, the sparse LU for systems of 1,000 equations or more and the
dense one below. \manref{cge_solve} takes \opt{sp}{arse} or \opt{den}{se} for one
solve.

\item[\optarg{}{mcp}{semismooth|activeset}] (from version 0.14.0) chooses the
method for the constraints of \manref{cge_constraint} and the inequalities of
\manref{cge_equation}: \cmd{semismooth}, every constraint in
Fischer--Burmeister form, solved at once by the semismooth Newton method (the
default), with the active-set method if it fails; or \cmd{activeset}, the
smoothed conditions, then the model solved in the regimes they give, a
constraint whose other condition fails changing regime
(\secref{sec:cons-solution}).
\end{optdesc}

\cgesec{Remarks and examples}

Both linear solvers give the same solution, to the solver's tolerance; they
differ in time and memory. The dense LU holds the Jacobian as a square matrix
and factorises it in time proportional to the cube of its size; the sparse LU
holds only its entries, chooses a pivot order once, which later solves of the
same model use again, and factorises each Jacobian in that order. Its first
solve of a model therefore costs more than its later ones. The full system of a
large model, with tens of thousands of equations, can only be solved with it.

The SAM is mapped and calibrated, and the tariffs removed:
\stlog{opt_setup}

\cgesubsec{Example 1: the sparse LU}

\stlog{opt_sparse}
The two solutions agree to the solver's tolerance, for the reduced and the
full system.

\cgesubsec{Example 2: the method for constraints}

Capacity limits on the three activities and quotas on the imports of two
commodities, with investment half as large again; four of the five bind.
Both methods find the same regimes and solution, the semismooth Newton method
in fewer iterations:
\stlog{opt_mcp}

\cgesec{Stored results}

\cmd{cge_options} stores the following in \cmd{r()}:
\begin{storedres}
\resgroup{Macros}
\res{r(linsolver)} & the linear solver: \cmd{dense}, \cmd{sparse} or \cmd{auto}\\
\res{r(mcp)} & the method for constraints: \cmd{semismooth} or \cmd{activeset}\\
\end{storedres}

\cgesec{Methods and formulas}

Section~\ref{sec:sparse} describes the sparse LU, and
section~\ref{sec:cons-solution} the methods for constraints.

\cgesec{Also see}

\begin{alsosee}
\seeentry{cge_solve}{Solve the model}
\seechapter{chap:solution}{The solution method}
\end{alsosee}
